\begin{glosario}
\begin{description}
  \item[Arquitectura modular:] Sistema de construcción que utiliza componentes prefabricados (módulos), que luego se ensamblan en el sitio final.
  \item[Capa de abstracción de hardware:] Componente de software que actúa como intermediario entre el sistema operativo o aplicación y el hardware físico, permitiendo interactuar con distintos dispositivos mediante una interfaz uniforme sin conocer sus detalles técnicos específicos.
  \item[Consistencia bidireccional:] Mecanismo de validación utilizado en diversos campos técnicos para asegurar que un proceso, dato o relación sea coherente tanto al analizarlo en una dirección como en la contraria.
  \item[Deuda técnica:] Costo futuro de retrabajo derivado de elegir soluciones rápidas y fáciles hoy, en lugar de un enfoque mejor diseñado pero más lento.
  \item[Eyetrackers:] Dispositivos tecnológicos que miden dónde mira una persona, la dirección de su mirada y el movimiento de sus ojos en tiempo real.
  \item[Interacción humano-información (HII):] Campo de investigación que estudia cómo las personas buscan, acceden, evalúan y usan información en distintos contextos y formatos.
  \item[MECIVA PS:] Metodología específica utilizada para la evaluación y validación de software científico, la cual permitió identificar los 67 errores estructurales en el sistema EDDIE actual.
  \item[Microservicios:] Arquitectura de software donde el sistema se divide en servicios independientes que se comunican mediante APIs, evaluada como una alternativa de solución para EDDIE por su alta escalabilidad.
  \item[Orquestador:] Sistema o herramienta que organiza, coordina y dispone diversos elementos para lograr un objetivo común, asegurando que funcionen de manera armoniosa.
  \item[Overhead de comunicación interproceso (IPC):] Tiempo y recursos adicionales de procesamiento que el sistema consume para permitir que módulos independientes intercambien información entre sí.
  \item[Plugins:] Aplicación informática pequeña que se integra en otro programa principal (como un navegador o gestor de contenido) para añadir funcionalidades específicas o mejorar su rendimiento sin modificar el código base.
  \item[Realidad aumentada:] Tecnología que superpone información digital, como imágenes, texto o datos, sobre el entorno físico real percibido por el usuario, sin reemplazarlo.
  \item[Science gateways:] Plataformas o portales que facilitan el acceso a recursos compartidos para la investigación científica, cuyas arquitecturas modulares sirven como referencia para la sostenibilidad de EDDIE.
  \item[SonarQube:] Plataforma de análisis estático de código diseñada para evaluar y mejorar continuamente la calidad, seguridad y mantenibilidad del software.
  \item[Visión computacional:] Área de la informática enfocada en permitir que las computadoras identifiquen y procesen objetos en imágenes o videos, utilizada en el proyecto para el reconocimiento de texto y gestos.
\end{description}
\end{glosario}
